% !Mode: : "TeX: UTF-8"
%%%(二)中英文摘要及关键词
%%%1．摘要是论文的内容不加注释和评论的简短陈述, 应做到不阅读全文就能获得与论文同等量的必要信息, 是一篇完整的短文, 可以独立使用、引用. 摘要应说明论文研究的目的、使用方法、结果和结论等. 
%%%2．中文摘要一般为1000字左右. 外文摘要为中文摘要的翻译, 所表述的内容与中文摘要一致. 外文摘要另起一页打印. 摘要中不用图、表、结构式公式及非公知公用的符号与术语. 
%%%3．关键词是为了文献标引工作而从论文中选取出来, 用以表示全论文主题内容信息的单词或术语, 一般为3—8个. 关键词字体加粗另起一行, 排在摘要的左下方, 按涉及的内容、从大到小排列. 多个关键词之间用分号隔开, 最后一个关键词不加标点, 应尽量用《汉语主题词表》等词表提供的规范词. 英文关键词、摘要内容应与本论文的中文关键词、摘要内容相对应. 
 

 %%%%%%% 中文摘要页: 

 %% 中文摘要内容
 \begin{cabstract}
《中国高考评价体系》是深化新时代高考内容改革的理论支撑和实践指南, 它指出学科素养是高考考查的关键内容之一. 与此同时, 逻辑推理素养作为高中数学六大核心素养之一, 不仅是数学严谨性的基本保证, 也是人们在数学活动中进行交流的基本思维品质. 因此, 如何基于高考评价体系进行逻辑推理素养的培养, 是目前一线教师重点关注的问题之一. 本文主要基于高考评价体系, 同时考虑《普通高中数学课程标准(2017年版2020年修订)》的要求, 以``圆锥曲线的方程"为切入点, 探索了高中生逻辑推理素养的培养策略, 在广东省汕尾市某高中进行了教学实验, 并检验了教学效果. 
\\
\indent
首先, 通过梳理相关文献, 界定了``推理与逻辑推理"和``逻辑推理素养"这两个核心概念. 然后, 分析了高考评价体系在数学学科的研究现状、逻辑推理素养的研究现状以及逻辑推理素养下圆锥曲线的教学现状. 发现虽然有较多学者关注逻辑推理素养的培养, 但关于考试与教学融合的研究工作尚不多见, 并且圆锥曲线与逻辑推理素养结合的研究较少.
\\
\indent
% 注: 在摘要环境中, 必须用上述两行来产生缩进的新段落
其次, 基于前面的文献研究, 结合高考评价体系, 对逻辑推理素养的水平划分框架与``四翼"的内涵进行了解读, 将逻辑推理素养与``四翼"均划分为三个水平等级, 水平越高意味着要求越高. 在此基础上, 选取四套新高考数学试卷, 以试卷中考查逻辑推理素养的试题为对象, 分析了其逻辑推理素养水平与``四翼"水平. 分析结果如下: (1) 试题对逻辑推理素养水平的考查主要集中在水平一与水平二, 且各卷在逻辑推理素养水平的命题情况基本一致;  (2) 试题对``四翼"的考查水平主要集中在基础性水平二与综合性水平二, 对应用性和创新性的考查较少. 
\\
\indent
然后, 依据《普通高中数学课程标准(2017年版2020年修订)》的要求编制了逻辑推理素养水平测试卷, 对广东省汕尾市某高中高二年级的学生进行了逻辑推理素养水平测试, 并以逻辑推理素养为主题, 对该校四位高二年级的数学教师进行了访谈. 利用SPSS分析了测试结果, 并总结了访谈结果, 得到以下结论: (1)大部分学生的逻辑推理素养水平处于水平一与水平二,极少学生能够达到水平三; 学生在合情推理方面的表现优于演绎推理; 学生的逻辑推理素养水平与学生的数学成绩呈正相关, 且男生与女生的逻辑推理素养水平不存在明显差异. (2)教师虽然无法准确地表述逻辑推理素养的内涵, 但能理解逻辑推理在数学学习中的重要性, 对圆锥曲线专题下培养学生逻辑推理素养的策略有一定思考; 虽然教师未对高考评价体系进行细致研究, 但对其核心内容都有一定的认识. 
\\
\indent
接着, 基于高考评价体系下的试卷分析结果, 结合高中生逻辑推理素养的水平调查与教师访谈, 提出了高考评价体系下高中生逻辑推理素养的培养策略: (1)注重直观与逻辑推理的辩证统一; (2)合理设置情境, 应对不同推理类型; (3)发展数学思维, 严谨逻辑表达; (4) 重视概念教学, 完善知识结构.
\\
\indent
最后, 依据制定的培养策略选取了实验班和对照班进行教学实验, 教学前后均对学生进行逻辑推理素养水平测试, 实验结果如下: (1)实验班学生的后测分数与前测分数几乎持平,对照班学生的后测分数相较于前测分数有所下降; (2)实验班在后测中达到水平三的人数比例高于前测; (3)在演绎推理方面, 实验班学生后测得分高于前测. 结果表明, 此次教学实验对学生逻辑推理素养的培养有一定的积极作用, 希望本文的研究能够为教学实践提供有益的参考. 
\end{cabstract}


 %% 中文关键词, 各关键词之间用逗号隔开, 最后一个关键词不加标点. 
 %% 关键词一般在3—8个之间. 
\begin{ckeywords}
逻辑推理素养, 高考评价体系, 圆锥曲线, 培养策略
\end{ckeywords}
 
 %%%%%%% 英文摘要页: 
 %%%%%%% 英文摘要内容和关键词应与本论文的中文关键词和摘要内容相对应. 
 
  %% 英文摘要内容
\begin{eabstract}	
The \textsl{`` China's College Entrance Examination Evaluation System"} is the theoretical support and practical guide for deepening the reform of the content of the new era college entrance examination. It points out that subject literacy is one of the key contents of the college entrance examination, and logical reasoning literacy, as one of the six core literacies in high school mathematics, is not only the basic guarantee of mathematical rigor but also the basic thinking quality for people to communicate in mathematical activities. Therefore, how to cultivate logical reasoning literacy based on the college entrance examination evaluation system is one of the key issues that frontline teachers currently focus on. Based on the college entrance examination evaluation system and considering the requirements of the \textsl{``National Curriculum Standard for General High School Mathematics (2017 Edition Revised in 2020)"}, this paper explores the cultivation strategy of logical reasoning literacy of high school students with the ``equation of conic curve" as the starting point, and conducts a teaching experiment in a high school in Shanwei City, Guangdong Province, and tests the teaching effect. 
\\
\indent
Firstly, by sorting out related literature, the core concepts of ``reasoning and logical reasoning" and ``logical reasoning literacy" were defined. Then, the research status of the college entrance examination evaluation system in mathematics, the research status of  logical reasoning literacy, and the teaching status of conic curve under logical reasoning literacy were analyzed. Although many scholars pay attention to the cultivation of logical reasoning literacy, there are few studies on the integration of testing and teaching, and there are few studies on the combination of conic curve and logical reasoning literacy. 
\\
\indent
Secondly, based on the previous literature research and the college entrance examination evaluation system, the level division framework and the ``four wings" connotation of logical reasoning literacy were interpreted, and both logical reasoning literacy and ``four wings" were divided into three levels, and the higher the level, the higher the requirements. On this basis, four sets of new college entrance examination mathematics papers were selected, and the logical reasoning literacy level and ``four wings" level of the questions on logical reasoning literacy in the papers were analyzed. The analysis results are as follows: (1) The examination of logical reasoning literacy level in the questions mainly focuses on level one and level two, and the proposition situation of the logical reasoning literacy level in each paper is basically the same; (2) The examination of ``four wings" level mainly focuses on foundational level two and comprehensive level two, and there is less examination of application and innovation. 
\\
\indent
Then, based on the requirements of the \textsl{``National Curriculum Standard for General High School Mathematics (2017 Edition Revised in 2020)"}, a test paper for logical reasoning literacy level was compiled, and a logical reasoning literacy level test was conducted on high school students in Grade 2 of a high school in Shanwei City, Guangdong Province, and four mathematics teachers in Grade 2 were interviewed on the topic of logical reasoning literacy. The test results were analyzed using SPSS, and the interview results were summarized. The following conclusions were drawn: (1) Most students' logical reasoning literacy level is at level one and level two, and very few students can reach level three; at the same time, their performance in conditional reasoning is better than deductive reasoning; there is a positive correlation between students' logical reasoning literacy level and their mathematical scores, and there is no significant difference in logical reasoning literacy level between male and female students. (2) Although teachers cannot accurately express the connotation of logical reasoning literacy, they can understand the importance of logical reasoning in mathematics learning and have some ideas on how to cultivate students' logical reasoning literacy under the topic of conic curve. Although teachers have not conducted detailed research on the college entrance examination evaluation system, they have a certain understanding of its core content. 
\\
\indent
Next, based on the analysis results of the college entrance examination evaluation system, the survey of high school students' logical reasoning literacy level, and the teacher interviews, the cultivation strategies for high school students' logical reasoning literacy under the college entrance examination evaluation system were proposed: (1) Emphasizing the dialectical unity of intuition and logical reasoning; (2) Reasonably setting up scenarios to deal with different types of reasoning; (3) Developing mathematical thinking and rigorous logical expression; (4) Attaching importance to concept teaching and improving knowledge structure. 
\\
\indent
Finally, based on the formulated cultivation strategies, an experimental class and a control class were selected for teaching experiments, and both classes were tested for their logical reasoning literacy level before and after the experiment. The results are as follows: (1) The post-test scores of the experimental class students are almost the same as the pre-test scores, while the post-test scores of the control class students have decreased compared to the pre-test scores; (2) The proportion of students who reach level three in the post-test of the experimental class is higher than that in the pre-test; (3) In deductive reasoning, the experimental class students score higher in the post-test than in the pre-test. The results show that this teaching experiment has a positive effect on cultivating students' logical reasoning literacy. It is hoped that this study can provide useful references for teaching practice. 
\end{eabstract}

 %% 英文关键词, 各关键词之间用逗号隔开. 外文关键词应与中文关键词相对应
\begin{ekeywords}
 logical reasoning literacy, college entrance examination evaluation system, conic curve, cultivation strategies.
\end{ekeywords} 
\makeabstract 